\section{商群的子群,商群的商群}

记号约定:$G$是拓扑群,$H\unlhd G$,$E$是拓扑空间,$\varphi:G\twoheadrightarrow G/H$是典范满射.

\begin{proposition}{拓扑群的第二同构定理}{}
	设$A\in\Sub\left(G\right)$,则典范双射$\phi\left(A\right)\to AH/H$是拓扑群同构.
\end{proposition}

\begin{remark}{}{}
	典范双射$A/\left(A\cap H\right)\to\phi\left(A\right)$是拓扑群同态,但一般不是拓扑群同构.
\end{remark}

\begin{proposition}{稠密子群的商群同构}{}
	设$G_{0}$是$G$的稠密子群,$H_{0}$是$G_{0}$的闭正规子群,$H=\closure\left(H_{0}\right)$,则典范双射$G_{0}/H_{0}\to\phi\left(G_{0}\right)$是拓扑群同构.
\end{proposition}

\begin{remark}{}{}
	\begin{enumerate}
		\item 设$G$连续作用于$E$,$K\unlhd G$.若$K$平凡地作用于$E$,则$G/K$连续地作用于$E$.
		\item 设$G$连续地作用于$E$,则$G$连续地作用于$E/H$.由$H$平凡作用于$E/H$可知,$G/H$连续地作用于$E/H$.
	\end{enumerate}
\end{remark}

\begin{proposition}{嵌套的轨道空间的拓扑群同构}{}
	设$G$连续地作用于$E$,则典范双射$E/G\to\left(E/H\right)/\left(G/H\right)$是同胚.
\end{proposition}

\begin{corollary}{拓扑群的第三同构定理}{}
	设$K\unlhd G$且$H\subseteq K$,则典范双射$G/K\to\left(G/H\right)/\left(K/H\right)$是拓扑群同构.
\end{corollary}














